% Prajñāpāramitāhṛdaya - Heart Sūtra
% Complete Multilingual Parallel Critical Edition
% Chinese · Sanskrit · Tibetan

\documentclass[10pt,a4paper,landscape]{article}

% ============================================================================
% PACKAGES
% ============================================================================

\usepackage[landscape,margin=1.2cm,bottom=2cm]{geometry}
\usepackage{fontspec}
\usepackage{xeCJK}
\usepackage{tabularx}
\usepackage{array}
\usepackage{booktabs}
\usepackage{xcolor}
\usepackage{longtable}
\usepackage{multicol}

% Per-page footnotes
\usepackage[perpage,symbol*]{footmisc}
\renewcommand{\thefootnote}{\arabic{footnote}}

% Page style
\usepackage{fancyhdr}
\pagestyle{fancy}
\fancyhf{}
\renewcommand{\headrulewidth}{0.4pt}
\fancyhead[L]{\small Prajñāpāramitāhṛdaya -- Complete Parallel Edition}
\fancyhead[R]{\small 般若波羅蜜多心經}
\fancyfoot[C]{\thepage}

% Hyperlinks
\usepackage{hyperref}
\hypersetup{colorlinks=true,linkcolor=blue}

% ============================================================================
% FONTS
% ============================================================================

\setmainfont{Noto Serif}
\setsansfont{Noto Sans}
\setCJKmainfont{Noto Serif CJK SC}
\setCJKsansfont{Noto Sans CJK SC}

\newfontfamily\devanagarifont{Noto Serif Devanagari}
\newcommand{\deva}[1]{{\devanagarifont #1}}

\newfontfamily\tibetanfont{Noto Serif Tibetan}
\newcommand{\tib}[1]{{\tibetanfont\small #1}}

% ============================================================================
% CUSTOM COMMANDS
% ============================================================================

\newcommand{\segnum}[1]{\textbf{[#1]}}
\newcommand{\rom}[1]{{\scriptsize\color{gray} #1}}
\newcommand{\eng}[1]{\multicolumn{3}{@{}p{0.96\textwidth}@{}}{\small\itshape #1} \\[0.3em]}
\newcommand{\var}[1]{\textsuperscript{\color{red}\textbf{#1}}}
\newcommand{\lemma}[1]{\textbf{#1}}
\newcommand{\reading}[1]{#1}
\newcommand{\wit}[1]{\textsc{#1}}
\newcommand{\depmark}[2]{{\scriptsize\color{blue}[#1$\to$#2]}}

% ============================================================================
% DOCUMENT
% ============================================================================

\begin{document}

% ----------------------------------------------------------------------------
% TITLE PAGE
% ----------------------------------------------------------------------------

\begin{center}
{\Huge\bfseries Prajñāpāramitāhṛdaya}\\[0.3em]
{\Huge 般若波羅蜜多心經}\\[0.3em]
{\huge\deva{प्रज्ञापारमिताहृदय}}\\[0.5em]
{\LARGE Heart Sūtra}\\[1em]
{\Large Complete Multilingual Parallel Critical Edition}\\[0.5em]
{\large Following Chinese Compositional Priority (Nattier 1992)}\\[2em]
\end{center}

\hrule
\vspace{1em}

\begin{multicols}{2}

\subsection*{Editorial Principles}

This edition recognizes that the Heart Sūtra was \textbf{composed in Chinese} (ca. 649--656 CE), with the Sanskrit being a back-translation and the Tibetan translated from Sanskrit.

\begin{itemize}
\item \textbf{Chinese T251} = Analytical base text (original composition)
\item \textbf{Sanskrit} = Back-translation, following Attwood's corrections
\item \textbf{Tibetan Toh 21} = Long recension, from Sanskrit
\item \textbf{Tibetan IOL Tib J 751} = Short recension (Dunhuang, c.\ 823 CE)
\end{itemize}

Key original finding: \textbf{T250 and T251 are parallel compositions} from Large Sūtra material, not revision stages of a single text.

\columnbreak

\subsection*{Witness Sigla}

\textit{Chinese:} \wit{T250} Dàmíngzhòujīng; \wit{T251} Xuanzang; \wit{Fangshan} 661 CE stele; \wit{T256} Tang transliteration; \wit{T257} Dānapāla 1005.

\textit{Sanskrit:} \wit{Ja} Hōryūji 7th-8th c.; \wit{Na} Nepal 11th c.; \wit{Nb Nc Nd Ne} Nepal 15th-19th c.

\textit{Tibetan:} \wit{D} Degé; \wit{Q} Peking; \wit{N} Narthang; \wit{H} Lhasa; \wit{J} Lithang; \wit{S} Stog; \wit{IOL} Dunhuang 823 CE (short recension).

\end{multicols}

\vspace{0.5em}
\hrule
\vspace{0.5em}

% ----------------------------------------------------------------------------
% SYNOPTIC TEXT
% ----------------------------------------------------------------------------

\footnotesize

\begin{longtable}{|>{\raggedright\arraybackslash}p{0.28\textwidth}|>{\raggedright\arraybackslash}p{0.34\textwidth}|>{\raggedright\arraybackslash}p{0.32\textwidth}|}
\hline
\textbf{Chinese (T251)} & \textbf{Sanskrit} & \textbf{Tibetan (D)} \\
\hline
\hline
\endfirsthead

\hline
\textbf{Chinese (T251)} & \textbf{Sanskrit} & \textbf{Tibetan (D)} \\
\hline
\hline
\endhead

\hline
\multicolumn{3}{r}{\small\itshape continued...}
\endfoot

\hline
\endlastfoot

% ============================================================================
% SEGMENT 1
% ============================================================================
\segnum{1} 觀自在\var{1}菩薩，行深般若波羅蜜多時，照見五蘊皆空，度一切苦厄\var{2}。
\newline\rom{Guānzìzài púsà, xíng shēn bōrě bōluómìduō shí, zhàojiàn wǔyùn jiē kōng, dù yīqiè kǔ è.}
&
ārya-avalokiteśvaro\var{3} bodhisattvo gambhīrāṃ prajñāpāramitā-caryāṃ caramāṇo vyavalokayati sma\var{4}. pañca-skandhāṃs\var{5} tāṃś ca svabhāva-śūnyān paśyati sma.
&
\tib{འཕགས་པ་སྤྱན་རས་གཟིགས་དབང་ཕྱུག་བྱང་ཆུབ་སེམས་དཔའ། ཤེས་རབ་ཀྱི་ཕ་རོལ་ཏུ་ཕྱིན་པ་ཟབ་མོའི་སྤྱོད་པ་སྤྱད་པ།}
\\
\hline
\eng{When Avalokiteśvara Bodhisattva was practicing the profound prajñāpāramitā, he illuminated and saw that the five aggregates are all empty, and crossed beyond all suffering.}
\hline
\hline

% ============================================================================
% SEGMENT 2
% ============================================================================
\segnum{2} 舍利子，色不異空，空不異色；色即是空，空即是色。受想行識亦復如是。
\newline\rom{Shèlìzǐ, sè bù yì kōng, kōng bù yì sè; sè jí shì kōng, kōng jí shì sè. Shòu xiǎng xíng shí yì fù rúshì.}
&
iha Śāriputra rūpaṃ śūnyatā śūnyataiva rūpam. rūpān na pṛthak śūnyatā śūnyatāyā na pṛthag rūpam. yad rūpaṃ sā śūnyatā yā śūnyatā tad rūpam\var{6}. evam eva vedanā-saṃjñā-saṃskāra-vijñānāni.
&
\tib{གཟུགས་སྟོང་པའོ། སྟོང་པ་ཉིད་གཟུགས་སོ། གཟུགས་ལས་སྟོང་པ་ཉིད་གཞན་མ་ཡིན། སྟོང་པ་ཉིད་ལས་ཀྱང་གཟུགས་གཞན་མ་ཡིན་ནོ།}
\\
\hline
\eng{Śāriputra, form is not different from emptiness, emptiness is not different from form. Form is emptiness, emptiness is form. Feeling, perception, volition, and consciousness are also like this.}
\hline
\hline

% ============================================================================
% SEGMENT 3
% ============================================================================
\segnum{3} 舍利子，是諸法空相，不生不滅，不垢不淨，不增不減。
\newline\rom{Shèlìzǐ, shì zhūfǎ kōng xiāng, bù shēng bù miè, bù gòu bù jìng, bù zēng bù jiǎn.}
&
iha Śāriputra sarva-dharmāḥ śūnyatā-lakṣaṇā. anutpannā aniruddhā. amalā avimalā. anūnā aparipūrṇāḥ.
&
\tib{དེ་ལྟར་ཆོས་ཐམས་ཅད་སྟོང་པ་ཉིད་དེ། མཚན་ཉིད་མེད་པ། མ་སྐྱེས་པ། མ་འགགས་པ། དྲི་མ་མེད་པ། དྲི་མ་དང་བྲལ་བ་མེད་པ། བྲི་བ་མེད་པ། གང་བ་མེད་པའོ།}
\\
\hline
\eng{Śāriputra, the emptiness-characteristic of all dharmas is: not arising, not ceasing; not defiled, not pure; not increasing, not decreasing.}
\hline
\hline

% ============================================================================
% SEGMENT 4
% ============================================================================
\segnum{4} 是故空中，無色，無受想行識，無眼耳鼻舌身意，無色聲香味觸法，
\newline\rom{Shìgù kōng zhōng, wú sè, wú shòu xiǎng xíng shí, wú yǎn ěr bí shé shēn yì, wú sè shēng xiāng wèi chù fǎ,}
&
tasmāc Chāriputra śūnyatāyāṃ na rūpaṃ na vedanā na saṃjñā na saṃskārāḥ na vijñānam. na cakṣuḥ-śrotra-ghrāṇa-jihvā-kāya-manāṃsi. na rūpa-śabda-gandha-rasa-spraṣṭavya-dharmāḥ.
&
\tib{དེ་ལྟ་བས་ན་སྟོང་པ་ཉིད་ལ་གཟུགས་མེད། ཚོར་བ་མེད། འདུ་ཤེས་མེད། འདུ་བྱེད་རྣམས་མེད། རྣམ་པར་ཤེས་པ་མེད། མིག་མེད། རྣ་བ་མེད། སྣ་མེད། ལྕེ་མེད། ལུས་མེད། ཡིད་མེད།}
\\
\hline
\eng{Therefore, in emptiness there is no form, no feeling, perception, volition, or consciousness; no eye, ear, nose, tongue, body, or mind; no form, sound, smell, taste, touch, or mental objects;}
\hline
\hline

% ============================================================================
% SEGMENT 5
% ============================================================================
\segnum{5} 無眼界，乃至無意識界。無無明，亦無無明盡\var{7}，乃至無老死，亦無老死盡。
\newline\rom{wú yǎn jiè, nǎizhì wú yìshí jiè. Wú wúmíng, yì wú wúmíng jìn, nǎizhì wú lǎosǐ, yì wú lǎosǐ jìn.}
&
na cakṣur-dhātur yāvan na manovijñāna-dhātuḥ. na-avidyā na-avidyā-kṣayo\var{8} yāvan na jarā-maraṇaṃ na jarā-maraṇa-kṣayaḥ.
&
\tib{མིག་གི་ཁམས་མེད་པ་ནས་ཡིད་ཀྱི་རྣམ་པར་ཤེས་པའི་ཁམས་ཀྱི་བར་དུ་ཡང་མེད་དོ། མ་རིག་པ་མེད། མ་རིག་པ་ཟད་པ་མེད་པ་ནས། རྒ་ཤི་མེད། རྒ་ཤི་ཟད་པའི་བར་དུ་ཡང་མེད་དོ།}
\\
\hline
\eng{no eye-element, up to and including no mind-consciousness element. No ignorance and no ending of ignorance, up to and including no old age and death, and no ending of old age and death.}
\hline
\hline

% ============================================================================
% SEGMENT 6
% ============================================================================
\segnum{6} 無苦集滅道，無智亦無得\var{9}，以無所得故。
\newline\rom{Wú kǔ jí miè dào, wú zhì yì wú dé, yǐ wú suǒ dé gù.}
&
na duḥkha-samudaya-nirodha-mārgāḥ. na jñānaṃ na prāptiḥ\var{10}.
&
\tib{སྡུག་བསྔལ་བ་དང་། ཀུན་འབྱུང་བ་དང་། འགོག་པ་དང་། ལམ་མེད། ཡེ་ཤེས་མེད། ཐོབ་པ་མེད། མ་ཐོབ་པ་ཡང་མེད་དོ།}
\\
\hline
\eng{No suffering, origination, cessation, or path. No wisdom and no attainment---because there is nothing to attain.}
\hline
\hline

% ============================================================================
% SEGMENT 7
% ============================================================================
\segnum{7} 菩提薩埵，依般若波羅蜜多故。心無罣礙，無罣礙故。無有恐怖，遠離顚倒夢想，究竟涅槃。
\newline\rom{Pútísàduǒ, yī bōrě bōluómìduō gù. Xīn wú guà'ài, wú guà'ài gù. Wú yǒu kǒngbù, yuǎnlí diāndǎo mèngxiǎng, jiūjìng nièpán.}
&
tasmāc Chāriputra aprāptitvād bodhisattvasya prajñāpāramitām āśritya viharaty acittāvaraṇaḥ. cittāvaraṇa-nāstitvād atrastro viparyāsa-atikrānto niṣṭhā-nirvāṇaḥ\var{11}.
&
\tib{བྱང་ཆུབ་སེམས་དཔའ་རྣམས་ཐོབ་པ་མེད་པའི་ཕྱིར། ཤེས་རབ་ཀྱི་ཕ་རོལ་ཏུ་ཕྱིན་པ་ལ་བརྟེན་ཅིང་གནས་ཏེ། སེམས་ལ་སྒྲིབ་པ་མེད་པས་འཇིགས་པ་མེད་དེ། ཕྱིན་ཅི་ལོག་ལས་ཤིན་ཏུ་འདས་ནས་མྱ་ངན་ལས་འདས་པའི་མཐར་ཕྱིན་ཏོ།}
\\
\hline
\eng{Because bodhisattvas rely on prajñāpāramitā, their minds are without obstruction. Because they are without obstruction, they have no fear. Far removed from inverted views and dreams, they attain ultimate nirvāṇa.}
\hline
\hline

% ============================================================================
% SEGMENT 8
% ============================================================================
\segnum{8} 三世諸佛，依般若波羅蜜多故，得阿耨多羅三藐三菩提。
\newline\rom{Sānshì zhū fó, yī bōrě bōluómìduō gù, dé ānòuduōluó sānmiǎo sānpútí.}
&
tryadhva-vyavasthitāḥ sarva-buddhāḥ prajñāpāramitām āśritya anuttarāṃ samyak-saṃbodhim abhisaṃbuddhāḥ.
&
\tib{དུས་གསུམ་དུ་རྣམ་པར་བཞུགས་པའི་སངས་རྒྱས་ཐམས་ཅད་ཀྱང་ཤེས་རབ་ཀྱི་ཕ་རོལ་ཏུ་ཕྱིན་པ་ལ་བརྟེན་ནས་བླ་ན་མེད་པ་ཡང་དག་པར་རྫོགས་པའི་བྱང་ཆུབ་ཏུ་མངོན་པར་རྫོགས་པར་སངས་རྒྱས་སོ།}
\\
\hline
\eng{All buddhas of the three times, relying on prajñāpāramitā, attain anuttarā-samyak-saṃbodhi (unsurpassed complete perfect awakening).}
\hline
\hline

% ============================================================================
% SEGMENT 9
% ============================================================================
\segnum{9} 故知般若波羅蜜多，是大神咒，是大明咒，是無上咒，是無等等咒\var{12}，能除一切苦，眞實不虛故。
\newline\rom{Gùzhī bōrě bōluómìduō, shì dà shén zhòu, shì dà míng zhòu, shì wúshàng zhòu, shì wú děngděng zhòu, néng chú yīqiè kǔ, zhēnshí bù xū gù.}
&
tasmāj jñātavyam: prajñāpāramitā mahā-mantro mahā-vidyā-mantro 'nuttara-mantro 'samasama-mantraḥ. sarva-duḥkha-praśamanaḥ. satyam amithyatvāt.
&
\tib{དེ་ལྟ་བས་ན་ཤེས་རབ་ཀྱི་ཕ་རོལ་ཏུ་ཕྱིན་པའི་སྔགས། རིག་པ་ཆེན་པོའི་སྔགས། བླ་ན་མེད་པའི་སྔགས། མི་མཉམ་པ་དང་མཉམ་པའི་སྔགས། སྡུག་བསྔལ་ཐམས་ཅད་རབ་ཏུ་ཞི་བར་བྱེད་པའི་སྔགས། མི་རྫུན་པས་ན་བདེན་པར་ཤེས་པར་བྱ་སྟེ།}
\\
\hline
\eng{Therefore know that prajñāpāramitā is the great mantra, the mantra of great illumination, the unsurpassed mantra, the unequalled mantra, which removes all suffering. It is true, not false.}
\hline
\hline

% ============================================================================
% SEGMENT 10: MANTRA
% ============================================================================
\segnum{10} 說般若波羅蜜多咒，卽說咒曰：揭帝揭帝，波羅揭帝，波羅僧揭帝，菩提僧莎訶\var{13}。
\newline\rom{Shuō bōrě bōluómìduō zhòu, jí shuō zhòu yuē: Jiēdì jiēdì, bōluó jiēdì, bōluó sēng jiēdì, pútí sēng suōhē.}
&
prajñāpāramitāyām ukto mantraḥ. tadyathā: [oṃ]\var{14} gate gate pāragate pārasaṃgate bodhi svāhā.
&
\tib{ཤེས་རབ་ཀྱི་ཕ་རོལ་ཏུ་ཕྱིན་པའི་སྔགས་སྨྲས་པ། ཏདྱ་ཐཱ། ཨོཾ་ག་ཏེ་ག་ཏེ་པཱ་ར་ག་ཏེ་པཱ་ར་སཾ་ག་ཏེ་བོ་དྷི་སྭཱ་ཧཱ།}
\\
\hline
\eng{The prajñāpāramitā mantra is proclaimed: Gate gate pāragate pārasaṃgate bodhi svāhā. (Gone, gone, gone beyond, gone completely beyond---awakening! Svāhā!)}
\hline

\end{longtable}

\newpage

% ============================================================================
% CRITICAL APPARATUS
% ============================================================================

\section*{Critical Apparatus}

\footnotesize

\begin{multicols}{2}

\var{1}~\lemma{觀自在}] \reading{觀世音} \wit{T250 T252}. 觀自在 reflects Xuanzang's translation of \textit{avalokita-īśvara}; 觀世音 reflects older \textit{avalokita-svara}. T250's pre-Xuanzang terminology confirms independent composition.

\var{2}~\lemma{度一切苦厄}] No Sanskrit or Tibetan parallel. Unique Chinese addition.

\var{3}~\lemma{avalokiteśvaro}] \reading{avalokita-} \wit{Ja}. Chinese 觀自在 supports \textit{-īśvara}. \depmark{zh}{sa}

\var{4}~\lemma{vyavalokayati sma}] Unusual imperfect. May reflect Chinese 時. \depmark{zh}{sa}

\var{5}~\lemma{pañca-skandhāṃs}] Accusative (correct). Conze erroneously prints nominative \reading{pañca-skandhāḥ}. Attwood (2015) corrects.

\var{6}~\lemma{yad rūpaṃ sā śūnyatā...}] This line has \textbf{no Chinese parallel}. Added in Sanskrit to clarify form-emptiness equation.

\var{7}~\lemma{盡} \textit{jìn}] = Sanskrit \textit{kṣaya}. Key back-translation evidence. See note 8.

\var{8}~\lemma{avidyā-kṣayo}] Uses \textit{kṣaya} where Large PP has \textit{nirodha}. ``Smoking gun'' (Nattier 1992:175--8). Chinese 盡 can mean either; back-translator chose \textit{kṣaya}. \depmark{zh}{sa} Tibetan \tib{ཟད་པ} \textit{zad pa} preserves error.

\var{9}~\lemma{無智亦無得}] T251, Fangshan, Ja lack \textit{na-aprāptiḥ}. Interpolation in later witnesses.

\var{10}~\lemma{na prāptiḥ}] Late witnesses add \reading{na-aprāptiḥ} \wit{T256 Nb-e}. Tibetan \tib{མ་ཐོབ་པ} preserves this interpolation.

\var{11}~\lemma{niṣṭhā-nirvāṇaḥ}] Conze prints \reading{-prāptaḥ}. Attwood (2018) questions compound.

\var{12}~\lemma{四咒 (four epithets)}] T251 has 4 epithets; T250 has 3 \textbf{exactly matching} T223 (Large Sūtra). Sanskrit follows T251's 4-epithet pattern. T250 preserves original.

\var{13}~\lemma{Mantra w/o oṃ}] T251, Fangshan, Ja begin directly with \textit{gate}. No oṃ.

\var{14}~\lemma{[oṃ]}] Bracketed: absent from early witnesses. Present in \wit{Nb Ne T257} and all Tibetan. Tantric addition (9th--10th c.).

\end{multicols}

\vspace{1em}
\hrule
\vspace{0.5em}

\section*{Stemmatic Summary}

\begin{center}
\begin{tabular}{lll}
\textbf{Stage} & \textbf{Date} & \textbf{Key Features} \\
\hline
α (T223 Large Sūtra) & 404 CE & Source passages (Kumārajīva) \\
T250/T251 (Chinese) & 5th--7th c. & \textbf{Parallel} compositions from α; T250 = 3 epithets, T251 = 4 epithets \\
σ (Sanskrit) & mid-7th c. & Back-translation from T251; \textit{kṣaya} error; \textit{yad rūpaṃ} addition \\
σ-L (Long) & 8th c. & Frame narrative added \\
Tibetan (IOL) & c.\ 823 CE & Dunhuang; short recension from σ; earliest datable Tibetan witness \\
Tibetan (Kangyur) & 8th--9th c. & From σ-L; preserves oṃ, \textit{na-aprāptiḥ}, \textit{zad pa} \\
τ (Tantric) & 9th--10th c. & oṃ in mantra; elaborate maṅgala \\
\end{tabular}
\end{center}

\newpage

% ----------------------------------------------------------------------------
% WORD-BY-WORD GLOSSARY
% ----------------------------------------------------------------------------

\section*{Word-by-Word Glossary}

Key terms glossed across all three traditions. Chinese is glossed at the phrase level; Sanskrit at the compound-analysis level; Tibetan at the word-group level.

\footnotesize

\subsection*{Core Vocabulary}

\begin{longtable}{p{3cm} p{3cm} p{3.5cm} p{3.5cm}}
\textbf{English} & \textbf{Chinese} & \textbf{Sanskrit} & \textbf{Tibetan} \\
\hline
\endhead
Avalokiteśvara & 觀自在 guānzìzài & avalokiteśvara & spyan ras gzigs dbang phyug \\
bodhisattva & 菩薩 púsà & bodhisattva & byang chub sems dpa' \\
prajñāpāramitā & 般若波羅蜜多 bōrě bōluómìduō & prajñāpāramitā & shes rab kyi pha rol tu phyin pa \\
five aggregates & 五蘊 wǔyùn & pañca-skandha & phung po lnga \\
form & 色 sè & rūpa & gzugs \\
emptiness & 空 kōng & śūnyatā & stong pa nyid \\
feeling & 受 shòu & vedanā & tshor ba \\
perception & 想 xiǎng & saṃjñā & 'du shes \\
formation & 行 xíng & saṃskāra & 'du byed \\
consciousness & 識 shí & vijñāna & rnam par shes pa \\
ignorance & 無明 wúmíng & avidyā & ma rig pa \\
ending/exhaustion & 盡 jìn & kṣaya ($\neq$ nirodha) & zad pa \\
old age \& death & 老死 lǎosǐ & jarā-maraṇa & rga shi \\
suffering & 苦 kǔ & duḥkha & sdug bsngal \\
path & 道 dào & mārga & lam \\
wisdom & 智 zhì & jñāna & ye shes \\
attainment & 得 dé & prāpti & thob pa \\
obstruction & 罣礙 guà'ài & āvaraṇa & sgrib pa \\
fear & 恐怖 kǒngbù & trasta & 'jigs pa \\
nirvāṇa & 涅槃 nièpán & nirvāṇa & mya ngan las 'das pa \\
mantra & 咒 zhòu & mantra & sngags \\
knowledge & 明 míng & vidyā & rig pa \\
gone & 揭帝 jiēdì & gate & ga te \\
awakening & 菩提 pútí & bodhi & byang chub \\
\end{longtable}

\normalsize

\vspace{1em}
\begin{center}
{\small\itshape Complete Parallel Critical Edition} $\cdot$ Prajñāpāramitāhṛdaya $\cdot$ Following Nattier (1992) and Attwood (2015--2021)
\end{center}

\end{document}
